
\documentclass[11pt]{article}
\usepackage[utf8]{inputenc}
\usepackage[spanish]{babel}
\usepackage{amsmath}
\usepackage{geometry}
\geometry{margin=1in}

\title{Documentación del Notebook SARIMA}
\date{}

\begin{document}
\maketitle

\section{Introducción}
Los modelos SARIMA extienden ARIMA incorporando componentes estacionales.

\section{Modelo General}
\[
SARIMA(p,d,q)(P,D,Q,s)
\]

donde:
\begin{itemize}
\item p: orden autorregresivo
\item d: diferenciación regular
\item q: media móvil
\item P: componente AR estacional
\item D: diferenciación estacional
\item Q: componente MA estacional
\item s: periodicidad estacional
\end{itemize}

\section{Herramientas de Identificación}
ACF y PACF permiten identificar componentes AR y MA.

\section{Selección del Modelo}
Se utilizan criterios AIC y BIC para comparar especificaciones.

\section{Diagnóstico}
Los residuos deben comportarse como ruido blanco.

\end{document}
